% =====================================================================
% NATHAN ROBERT - ONE-PAGE CV (Overleaf / pdflatex)
% Reference: classic minimalist, single-column resume.
% HOW TO PERSONALIZE: complete the fields below, including LinkedIn.
% Once LinkedInURL is filled in, a SCANNABLE QR CODE is generated.
% No additional images are required.
% =====================================================================
\documentclass[10pt,a4paper]{extarticle}
\usepackage[T1]{fontenc}
\usepackage[utf8]{inputenc}
\usepackage[french]{babel}
\usepackage[scaled=0.96]{helvet}
\renewcommand{\familydefault}{\sfdefault}
\usepackage[margin=17.5mm,top=17mm,bottom=16mm]{geometry}
\usepackage{microtype}
\usepackage{enumitem}
\usepackage{etoolbox}
\usepackage{hyperref}
\usepackage{qrcode}
\usepackage{fontawesome5}
\usepackage{xcolor}
\usepackage{tabularx}
\usepackage{array}
\hypersetup{colorlinks=true,urlcolor=black,pdfauthor={Nathan Robert},pdftitle={CV - Nathan Robert}}
\setlength{\parindent}{0pt}
\setlength{\parskip}{0pt}
\linespread{1.015}
\pagestyle{empty}
\definecolor{CVGray}{gray}{0.33}

% -------- EDIT THESE FOUR FIELDS --------------------------------------
\newcommand{\CVEmail}{} % e.g. name@email.com
\newcommand{\CVPhone}{} % e.g. +33 6 ...
\newcommand{\LinkedInURL}{https://www.linkedin.com/in/nathan-robert-b9423a320/} % e.g. https://www.linkedin.com/in/your-profile/
\newcommand{\CVGitHub}{} % e.g. https://github.com/your-username
% ---------------------------------------------------------------------

\newcommand{\contactitem}[2]{%
  \ifdefempty{#2}{}{#1\,#2\par}%
}
\newcommand{\renderemail}{%
  \ifdefempty{\CVEmail}{}{\faEnvelope\,\href{mailto:\CVEmail}{\CVEmail}\par}%
}
\newcommand{\renderphone}{%
  \ifdefempty{\CVPhone}{}{\faPhone\,\CVPhone\par}%
}
\newcommand{\renderlinkedin}{%
  \ifdefempty{\LinkedInURL}{}{\faLinkedin\,\href{\LinkedInURL}{Profil LinkedIn}\par}%
}
\newcommand{\rendergithub}{%
  \ifdefempty{\CVGitHub}{}{\faGithub\,\href{\CVGitHub}{Profil GitHub}\par}%
}
\newcommand{\renderqr}{%
  \ifdefempty{\LinkedInURL}{%
    \fbox{\begin{minipage}[c][17mm][c]{17mm}\centering
      {\fontsize{9}{10}\selectfont\bfseries QR}\\[2pt]
      {\fontsize{6.2}{7}\selectfont LINKEDIN}
    \end{minipage}}%
  }{\href{\LinkedInURL}{\qrcode[height=17mm]{\LinkedInURL}}}%
}
\newcommand{\cvsection}[1]{%
  \vspace{11pt}
  {\fontsize{11.1}{12.7}\selectfont\bfseries\MakeUppercase{#1}}\par
  \vspace{2.5pt}\hrule height .50pt\vspace{5pt}
}
\newcommand{\cvrole}[3]{%
  {\fontsize{10}{12}\selectfont\bfseries #1}\hfill
  {\fontsize{9.1}{11}\selectfont #2}\par
  {\fontsize{9.0}{11}\selectfont\color{CVGray}#3}\par
  \vspace{2pt}
}
\newenvironment{cvbullets}{%
  \begin{itemize}[leftmargin=13pt,label={\textbullet},topsep=2pt,partopsep=0pt,
    parsep=0pt,itemsep=1.3pt]
}{\end{itemize}}
\newcommand{\cvsmallgap}{\vspace{3.8pt}}

\begin{document}
\noindent
\begin{minipage}[c]{0.54\textwidth}
  {\fontsize{22.5}{26}\selectfont\bfseries Nathan ROBERT}\par
  \vspace{3.5pt}
  {\fontsize{10.4}{12.5}\selectfont IA \& Machine Learning $\vert$ NLP, Fine-Tuning de LLM}\par
\end{minipage}%
\hfill
\begin{minipage}[c]{0.435\textwidth}
\setlength{\tabcolsep}{0pt}
\begin{tabular}{@{}!{\vrule width .5pt}@{\hspace{9pt}}p{.64\linewidth}@{\hspace{5pt}}c@{}}
  \begin{minipage}[c]{\linewidth}
    {\fontsize{8.3}{10.6}\selectfont
    \faMapMarker*\, Toulouse, France\par
    \renderemail\renderphone\renderlinkedin\rendergithub}
  \end{minipage}
  & \renderqr
\end{tabular}
\end{minipage}
\vspace{5pt}


\cvsection{Profil}
Étudiant ingénieur en informatique à l'ENSEEIHT, orienté intelligence artificielle et traitement automatique du langage. Après un stage en fine-tuning de LLM open-source chez CPMS, j'y développe depuis octobre 2026 des solutions d'IA et de machine learning en freelance.

\cvsection{Compétences}
\textbf{IA / NLP :} Transformers, modèles de type BERT, classification documentaire, fine-tuning supervisé de LLM, LoRA, QLoRA, évaluation.\\[1.8pt]
\textbf{Bibliothèques :} PyTorch, Hugging Face (Transformers, Datasets, PEFT, TRL), Unsloth, bitsandbytes.\\[1.8pt]
\textbf{Programmation / outils :} Python, C, OCaml, Ada, SQL, Linux ; préparation de données, tokenisation, optimisation GPU.

\cvsection{Expérience professionnelle}
\cvrole{Freelance IA \& Machine Learning | CPMS}{Oct. 2026 -- Aujourd'hui}{Centre de Prévoyance Médico-Sociale \textbar\ Toulouse \textbar\ À distance}
\begin{cvbullets}
  \item Développement de solutions de classification de documents à l'aide du machine learning et de modèles NLP de type BERT.
  \item Adaptation de modèles de langage par fine-tuning, évaluation des performances et analyse des erreurs.
  \item Préparation et traitement de datasets, participation à l'intégration des modèles d'IA dans des applications.
\end{cvbullets}
\cvsmallgap
\cvrole{Stagiaire IA générative / Fine-Tuning de LLM | CPMS}{Juin -- Juil. 2026}{Centre de Prévoyance Médico-Sociale \textbar\ Paris \textbar\ Sur site}
\begin{cvbullets}
  \item Étude des architectures Transformer et adaptation de LLM open-source à un cas d'usage métier impliquant des documents médicaux sensibles et des enjeux RGPD.
  \item Constitution de datasets instruction--réponse, nettoyage, normalisation, tokenisation et gestion des chat templates.
  \item Expérimentations de fine-tuning supervisé avec LoRA / QLoRA et quantification 4/8 bits pour optimiser l'utilisation GPU.
  \item Réglage des hyperparamètres : learning rate, batch size, epochs, accumulation des gradients et warmup.
  \item Comparaison prompting / fine-tuning, analyses d'erreurs, itérations et prise en compte du surapprentissage, des hallucinations et de la généralisation.
\end{cvbullets}

\cvsection{Formation}
\cvrole{Diplôme d'ingénieur en informatique | ENSEEIHT -- Toulouse INP}{2025 -- 2028}{Toulouse \textbar\ Formation en cours}
\vspace{2pt}
\cvrole{Classes préparatoires (CPGE MP2I / MPI) | Lycée Janson-de-Sailly}{2022 -- 2025}{Paris}

\cvsection{Projets techniques}
\textbf{2025} \textbar\ Conception d'un modèle de prédiction de l'expansion urbaine.\\[1.7pt]
\textbf{2024} \textbar\ Développement d'un simulateur de match de football en 2D.\\[1.7pt]
\textbf{2023} \textbar\ Développement d'un jeu vidéo.

\cvsection{Langues \& engagements}
\textbf{Langues :} Français (natif), anglais (C1), espagnol (intermédiaire).\\[1.7pt]
\textbf{Bénévolat :} Service et distribution de repas aux Restaurants du Cœur (2016--2022).\\[1.7pt]
\textbf{Centres d'intérêt :} Force athlétique, football, programmation.
% Expérience plus ancienne du CV d'origine, non affichée :
% stage de découverte informatique chez Alan (février 2019).


\end{document}
